\chapter{Continuité, valeurs intermédiaires et équations}
\section*{Exercice 13.1 -- Existence et unicité}
[Méthode] Montrer que $x^3+2x-2=0$ a une unique solution réelle, située dans $]0,1[$.

\section*{Exercice 13.2 -- Dichotomie}
[Base] Combien de dichotomies sur $[0,1]$ garantissent que le milieu approche une racine à moins de $10^{-3}$ ?

\section*{Exercice 13.3 -- Prolongement}
[Méthode] Prolonger continûment $(e^x-1)/x$ en 0.

\section*{Exercice 13.4 -- TVI sans monotonie}
[Critique] L'équation $x^2=1/4$ sur $[-1,1]$ a-t-elle une solution unique ?

\section*{Exercice 13.5 -- Image d’une suite}
[Preuve] Si $u_n\to2$, calculer la limite de $\sqrt{1+u_n^2}$.

\section*{Exercice 13.6 -- Additivité continue}
[Approfondissement] Si $f$ est continue et additive avec $f(1)=3$, déterminer $f(\sqrt2)$.
